\chapter{2008年福建卷（理）}

\section{单选题}

\begin{problem}
若复数 \((a^2 - 3a + 2) + (a - 1)\i\) 是纯虚数，则实数 \(a\) 的值为（\quad）

\choices
    {\(1\)}
    {\(2\)}
    {\(1\) 或 \(2\)}
    {\(-1\)}
\begin{answer}
B
\end{answer}
\end{problem}

\begin{problem}
设集合 \(A = \left\{x \mid \frac{x}{x - 1} < 0\right\}\)，\(B = \{x \mid 0 < x < 3\}\)，那么“\(m \in A\)”是“\(m \in B\)”的（\quad）

\choices
    {充分而不必要条件}
    {必要而不充分条件}
    {充要条件}
    {既不充分也不必要条件}
\begin{answer}
A
\end{answer}
\end{problem}

\begin{problem}
设 \(\{a_n\}\) 是公比为正数的等比数列，若 \(a_1 = 7\)，\(a_5 = 16\)，则数列 \(\{a_n\}\) 前 7 项的和为（\quad）

\choices
    {\(63\)}
    {\(64\)}
    {\(127\)}
    {\(128\)}
\end{problem}

\begin{solution}
设等比数列的公比为 \(q\)，则 \(q>0\). 由 \(a_5=a_1q^4\)，得

\[
q^4=\frac{a_5}{a_1}=\frac{16}{7}.
\]

又因为

\[
\left(\frac{6}{5}\right)^4=\frac{1296}{625}<\frac{16}{7}<\frac{625}{256}=\left(\frac{5}{4}\right)^4,
\]

所以

\[
\frac{6}{5}<q<\frac{5}{4}.
\]

又因为 \(q>1\)，所以 \(S_7=7\left(1+q+q^2+\cdots+q^6\right)\) 随 \(q\) 严格增大，从而

\[
35\left[\left(\frac{6}{5}\right)^7-1\right]
<S_7
<28\left[\left(\frac{5}{4}\right)^7-1\right].
\]

其中

\[
35\left[\left(\frac{6}{5}\right)^7-1\right]=\frac{7063385}{78125}>64.
\]

并且

\[
28\left[\left(\frac{5}{4}\right)^7-1\right]=\frac{1728748}{16384}<127.
\]

故 \(64<S_7<127\)，而选项只有 \(63\)、\(64\)、\(127\)、\(128\)，没有一个与 \(S_7\) 相等. 因此，按原试卷的题面，题设条件与选项不相容，无法确定选择项.
\end{solution}

\begin{problem}
函数 \( f(x) = x^3 + \sin x + 1 \)（\(x \in \R\)），若 \( f(a) = 2 \)，则 \( f(-a) \) 的值为（\quad）

\choices
    {\(3\)}
    {\(0\)}
    {\(-1\)}
    {\(-2\)}
\begin{answer}
B
\end{answer}
\end{problem}

\begin{problem}
某一批花生种子，如果每 \(1\) 粒发芽的概率为 \(\frac{4}{5}\)，那么播下 \(4\) 粒种子恰有 \(2\) 粒发芽的概率是（\quad）

\choices
    {\(\frac{16}{625}\)}
    {\(\frac{96}{625}\)}
    {\(\frac{192}{625}\)}
    {\(\frac{256}{625}\)}
\begin{answer}
B
\end{answer}
\end{problem}

\begin{problem}
如图，在长方体 \(ABCD - A_{1}B_{1}C_{1}D_{1}\) 中，\(AB = BC = 2\)，\(AA_{1} = 1\)，则 \(BC_{1}\) 与平面 \(BB_{1}D_{1}D\) 所成角的正弦值为（\quad）

\bitmapfigure[max width=0.42\linewidth]{img/2008/fujian_science/q6_fig1.png}

\choices
    {\(\frac{\sqrt{6}}{3}\)}
    {\(\frac{2\sqrt{6}}{5}\)}
    {\(\frac{\sqrt{15}}{5}\)}
    {\(\frac{\sqrt{10}}{5}\)}
\begin{answer}
D
\end{answer}
\end{problem}

\begin{problem}
某班级要从 \(4\) 名男生，\(2\) 名女生中选派 \(4\) 人参加某次社区服务，如果要求至少有 \(1\) 名女生，那么不同的选派方案种数为（\quad）

\choices
    {\(14\)}
    {\(24\)}
    {\(28\)}
    {\(48\)}
\begin{answer}
A
\end{answer}
\end{problem}

\begin{problem}
若实数 \(x\)， \(y\) 满足 \(\begin{cases}
x - y + 1 \leqslant 0,\\
x > 0,\\
\end{cases}\) 则 \(\frac{y}{x}\) 的取值范围是（\quad）

\choices
    {\((0,1)\)}
    {\((0,1]\)}
    {\((1, +\infty)\)}
    {\([1, +\infty)\)}
\begin{answer}
C
\end{answer}
\end{problem}

\begin{problem}
函数 \( f(x) = \cos x \)（\(x \in \R\)）的图象按向量 \((m, 0)\) 平移后，得到函数 \( y = -f'(x) \) 的图象，则 \( m \) 的值可以为（\quad）

\choices
    {\(\frac{\pi}{2}\)}
    {\(\pi\)}
    {\(-\pi\)}
    {\(-\frac{\pi}{2}\)}
\begin{answer}
A
\end{answer}
\end{problem}

\begin{problem}
在 \(\triangle ABC\) 中，角 \(A\)、\(B\)、\(C\) 的对边分别为 \(a\)、\(b\)、\(c\)，若 \((a^{2} + c^{2} - b^{2})\tan B = \sqrt{3}ac\)，则角 \(B\) 的值为（\quad）

\choices
    {\(\frac{\pi}{6}\)}
    {\( \frac{\pi}{3} \)}
    {\(\frac{\pi}{6}\) 或 \(\frac{5\pi}{6}\)}
    {\(\frac{\pi}{3}\) 或 \(\frac{2\pi}{3}\)}
\begin{answer}
D
\end{answer}
\end{problem}

\begin{problem}
双曲线 \(\frac{x^2}{a^2} - \frac{y^2}{b^2} = 1\)（\(a > 0\)，\(b > 0\)）的两个焦点为 \(F_1\)，\(F_2\)，若 \(P\) 为其上一点，且 \(|PF_1| = 2|PF_2|\)，则双曲线离心率的取值范围为（\quad）

\choices
    {\((1,3)\)}
    {\((1,3]\)}
    {\( (3, +\infty) \)}
    {\([3, +\infty)\)}
\begin{answer}
B
\end{answer}
\end{problem}

\begin{problem}
已知函数 \( y = f(x) \)，\( y = g(x) \) 的导函数的图象如下图，那么 \( y = f(x) \)，\( y = g(x) \) 的图象可能是（\quad）

\bitmapfigure[max width=0.88\linewidth]{img/2008/fujian_science/q12_fig1.png}

\choices
    {\choicebitmap[width=0.15\paperwidth]{img/2008/fujian_science/q12_optA.png}}
    {\choicebitmap[width=0.15\paperwidth]{img/2008/fujian_science/q12_optB.png}}
    {\choicebitmap[width=0.15\paperwidth]{img/2008/fujian_science/q12_optC.png}}
    {\choicebitmap[width=0.15\paperwidth]{img/2008/fujian_science/q12_optD.png}}
\begin{answer}
D
\end{answer}
\end{problem}

\section{填空题}

\begin{problem}
若 \((x - 2)^{5} = a_{5}x^{5} + a_{4}x^{4} + a_{3}x^{3} + a_{2}x^{2} + a_{1}x + a_{0}\)，则 \(a_1 + a_2 + a_3 + a_4 + a_5 =\fillinblank{}.\)
\begin{answer}
31
\end{answer}
\end{problem}

\begin{problem}
若直线 \(3x + 4y + m = 0\) 与圆 \(\begin{cases}
x = 1 + \cos \theta,\\
y = -2 + \sin \theta,
\end{cases}\)（\(\theta\) 为参数）没有公共点，则实数 \(m\) 的取值范围是\fillinblank{}.
\begin{answer}
\(m<0\) 或 \(m>10\)
\end{answer}
\end{problem}

\begin{problem}
若三棱锥的三个侧面两两垂直，且侧棱长均为 \( \sqrt{3} \)，则其外接球的表面积是\fillinblank{}.
\begin{answer}
\(9\pi\)
\end{answer}
\end{problem}

\begin{problem}
设 \(P\) 是一个数集，且至少含有两个数，若对任意 \(a\in P\)，\(b \in P\)，都有 \(a + b\)，\(a - b\)，\(ab\)，\(\frac{a}{b} \in P\)（除数 \(b \neq 0\)），则称 \(P\) 是一个数域. 例如有理数集 \(\mathbf{Q}\) 是数域；数集 \(F = \{a + b\sqrt{2} \mid a, b \in \mathbf{Q}\}\) 也是数域. 有下列命题：

\begin{circlelist}
\item 整数集是数域；
\item 若有理数集 \(\mathbf{Q} \subseteq M\)，则数集 \(M\) 必为数域；
\item 数域必为无限集；
\item 存在无穷多个数域.
\end{circlelist}

其中正确的命题的序号是\fillinblank{}.
\begin{answer}
\(3\)、\(4\)
\end{answer}
\end{problem}

\section{解答题}

\begin{problem}
已知向量 \(\bs{m} = (\sin A, \cos A)\)，\(\bs{n} = (\sqrt{3}, -1)\)，\(\bs{m} \cdot \bs{n} = 1\)，且 \(A\) 为锐角.
\begin{enumerate}
\item 求角 \(A\) 的大小；
\item 求函数 \( f(x) = \cos 2x + 4 \cos A \sin x \)（\(x \in \R\)）的值域.
\end{enumerate}
\end{problem}

\begin{solution}
\begin{enumerate}
\item 由 \(\bs{m} \cdot \bs{n} = 1\)，得
\[
\sqrt{3}\sin A-\cos A=1.
\]
注意到 \(\sqrt{3}\sin A-\cos A=2\sin\left(A-\frac{\pi}{6}\right)\)，故
\[
\sin\left(A-\frac{\pi}{6}\right)=\frac{1}{2}.
\]
因为 \(A\) 为锐角，所以 \(A-\frac{\pi}{6}\in\left(-\frac{\pi}{6},\frac{\pi}{3}\right)\)，在此范围内只有 \(A-\frac{\pi}{6}=\frac{\pi}{6}\) 满足条件，从而 \(A=\frac{\pi}{3}\).

\item 代入 \(\cos A=\frac{1}{2}\)，得
\[
f(x)=\cos 2x+2\sin x=1-2\sin^2x+2\sin x=\frac{3}{2}-2\left(\sin x-\frac{1}{2}\right)^2.
\]
令 \(t=\sin x\)，则 \(t\in[-1,1]\). 上式在 \(t=\frac{1}{2}\) 时取得最大值 \(\frac{3}{2}\)，在 \(t=-1\) 时取得最小值 \(-3\). 因此，函数 \(f(x)\) 的值域为 \(\left[-3,\frac{3}{2}\right]\).
\end{enumerate}
\end{solution}

\begin{problem}
如图，在四棱锥 \(P - ABCD\) 中，侧面 \(PAD \perp\) 底面 \(ABCD\)，侧棱 \(PA = PD = \sqrt{2}\)，底面 \(ABCD\) 为直角梯形，其中 \(BC \parallel AD\)，\(AB \perp AD\)，\(AD = 2AB = 2BC = 2\)，\(O\) 为 \(AD\) 中点.

\begin{enumerate}
\item 求证： \(PO \perp\) 平面 \(ABCD\)；
\item 求异面直线 \(PD\) 与 \(CD\) 所成角的大小；
\item 线段 \(AD\) 上是否存在点 \(Q\)，使得它到平面 \(PCD\) 的距离为 \(\frac{\sqrt{3}}{2}\). 若存在，求出 \(\frac{AQ}{QD}\) 的值. 若不存在，请说明理由.
\end{enumerate}

\bitmapfigure[max width=0.36\linewidth]{img/2008/fujian_science/q18_fig1.png}
\end{problem}

\begin{solution}
建立空间直角坐标系，使底面 \(ABCD\) 在 \(xOy\) 平面内，并取
\(A(0,0,0)\)，\(D(2,0,0)\)，\(B(0,1,0)\)，\(C(1,1,0)\).
则 \(O(1,0,0)\). 由 \(PAD\perp ABCD\)，平面 \(PAD\) 可取为 \(xOz\) 平面；又因为 \(PA=PD\)，且 \(O\) 是 \(AD\) 的中点，所以 \(PO\perp AD\). 于是 \(PO\perp\) 平面 \(ABCD\). 另外，\(AO=1\)，故
\[
PO=\sqrt{PA^2-AO^2}=\sqrt{2-1}=1.
\]
取 \(P(1,0,1)\).

\begin{enumerate}
\item 上述论证已证明 \(PO\perp\) 平面 \(ABCD\).

\item 取直线 \(PD\) 和 \(CD\) 的方向向量分别为
\(\overrightarrow{DP}=(-1,0,1)\) 和 \(\overrightarrow{DC}=(-1,1,0)\). 设异面直线所成的角为 \(\theta\)，则
\[
\cos\theta=\frac{\left|\overrightarrow{DP}\cdot\overrightarrow{DC}\right|}{|\overrightarrow{DP}|\,|\overrightarrow{DC}|}=\frac{1}{\sqrt{2}\cdot\sqrt{2}}=\frac{1}{2}.
\]
因此 \(\theta=\frac{\pi}{3}\).

\item 设 \(Q(q,0,0)\)，其中 \(0\leqslant q\leqslant2\). 由
\(\overrightarrow{PC}=(0,1,-1)\) 和 \(\overrightarrow{PD}=(1,0,-1)\)，
可取平面 \(PCD\) 的一个法向量为 \((1,1,1)\)，故平面 \(PCD\) 的方程为 \(x+y+z=2\). 于是
\[
d(Q,\text{平面 }PCD)=\frac{|q-2|}{\sqrt{3}}=\frac{2-q}{\sqrt{3}}.
\]
令其等于 \(\frac{\sqrt{3}}{2}\)，得 \(q=\frac{1}{2}\)，此时
\[
\frac{AQ}{QD}=\frac{q}{2-q}=\frac{1/2}{3/2}=\frac{1}{3}.
\]
所以存在这样的点 \(Q\)，且 \(\frac{AQ}{QD}=\frac{1}{3}\).
\end{enumerate}
\end{solution}

\begin{problem}
已知函数 \( f(x) = \frac{1}{3}x^3 + x^2 - 2 \).
\begin{enumerate}
\item 设 \(\{a_n\}\) 是正数组成的数列. 前 \(n\) 项和为 \(S_n\)，其中 \(a_{1} = 3\). 若点 \((a_{n}, a_{n+1}^{2} - 2a_{n+1})\)（\(n \in \N^*\)）在函数 \(y = f'(x)\) 的图象上，求证：点 \((n, S_{n})\) 也在 \(y = f'(x)\) 的图象上；
\item 求函数 \( f(x) \) 在区间 \( (a-1, a) \) 内的极值.
\end{enumerate}
\end{problem}

\begin{solution}
由 \(f'(x)=x^2+2x=x(x+2)\)，下面逐问求解.
\begin{enumerate}
\item 点 \((a_n,a_{n+1}^2-2a_{n+1})\) 在 \(y=f'(x)\) 的图象上，等价于
\[
a_{n+1}^2-2a_{n+1}=a_n^2+2a_n.
\]
整理得
\[
(a_{n+1}-1)^2=(a_n+1)^2.
\]
因此 \(a_{n+1}-1=a_n+1\) 或 \(a_{n+1}-1=-(a_n+1)\). 第二种情形给出 \(a_{n+1}=-a_n<0\)，与数列各项为正数矛盾，所以 \(a_{n+1}=a_n+2\). 结合 \(a_1=3\)，得
\[
a_n=3+2(n-1)=2n+1.
\]
从而
\[
S_n=\frac{n(a_1+a_n)}{2}=\frac{n(3+2n+1)}{2}=n^2+2n=f'(n).
\]
故点 \((n,S_n)\) 在 \(y=f'(x)\) 的图象上.

\item 因为 \(f'(x)=x(x+2)\)，所以 \(f(x)\) 在 \((-\infty,-2)\) 上递增，在 \((-2,0)\) 上递减，在 \((0,+\infty)\) 上递增. 又 \(f(-2)=-\frac{2}{3}\)，\(f(0)=-2\).
当 \(-2<a<-1\) 时，区间 \((a-1,a)\) 内含有 \(-2\) 而不含 \(0\)，故 \(f(x)\) 在 \(x=-2\) 处取得极大值 \(-\frac{2}{3}\). 当 \(0<a<1\) 时，区间 \((a-1,a)\) 内含有 \(0\) 而不含 \(-2\)，故 \(f(x)\) 在 \(x=0\) 处取得极小值 \(-2\). 其余情形下，区间 \((a-1,a)\) 不含 \(-2\) 或 \(0\)，函数在该区间内单调，因而没有极值.
\end{enumerate}
\end{solution}

\begin{problem}
某项考试按科目 \(A\)、科目 \(B\) 依次进行，只有当科目 \(A\) 成绩合格时，才可继续参加科目 \(B\) 的考试. 已知每个科目只允许有一次补考机会，两个科目成绩均合格方可获得证书. 现某人参加这项考试，科目 \(A\) 每次考试成绩合格的概率均为 \( \frac{2}{3} \)，科目 \(B\) 每次考试成绩合格的概率均为 \( \frac{1}{2} \). 假设各次考试成绩合格与否均互不影响.
\begin{enumerate}
\item 求他不需要补考就可获得证书的概率；
\item 在这项考试过程中，假设他不放弃所有的考试机会，记他参加考试的次数为 \(\xi\)，求 \(\xi\) 的数学期望 \(E(\xi)\).
\end{enumerate}
\end{problem}

\begin{solution}
\begin{enumerate}
\item 不需要补考就获得证书，必须第一次参加科目 \(A\) 和第一次参加科目 \(B\) 都合格，所以所求概率为
\[
\frac{2}{3}\times\frac{1}{2}=\frac{1}{3}.
\]

\item 科目 \(A\) 至少参加一次，且只有第一次不合格时才参加补考，因此参加科目 \(A\) 的次数的数学期望为
\[
1+\frac{1}{3}=\frac{4}{3}.
\]
科目 \(A\) 在两次机会内至少合格一次的概率为
\[
1-\left(\frac{1}{3}\right)^2=\frac{8}{9}.
\]
只有科目 \(A\) 合格后才参加科目 \(B\). 在已经参加科目 \(B\) 的条件下，第一次考试必然参加，第一次不合格的概率为 \(\frac{1}{2}\)，所以参加科目 \(B\) 的次数的条件数学期望为
\[
1+\frac{1}{2}=\frac{3}{2}.
\]
因此参加科目 \(B\) 的次数的数学期望为 \(\frac{8}{9}\times\frac{3}{2}=\frac{4}{3}\). 由期望的可加性，
\[
E(\xi)=\frac{4}{3}+\frac{4}{3}=\frac{8}{3}.
\]
\end{enumerate}
\end{solution}

\begin{problem}
如图，椭圆 \(\frac{x^2}{a^2} + \frac{y^2}{b^2} = 1\)（\(a > b > 0\)）的一个焦点是 \(F(1,0)\)，\(O\) 为坐标原点.

\begin{enumerate}
\item 已知椭圆短轴的两个三等分点与一个焦点构成正三角形，求椭圆的方程；
\item 设过点 \(F\) 的直线 \(l\) 交椭圆于 \(A\)、\(B\) 两点. 若直线 \(l\) 绕点 \(F\) 任意转动，恒有 \(|OA|^2 + |OB|^2 < |AB|^2\)，求 \(a\) 的取值范围.
\end{enumerate}

\bitmapfigure[max width=0.34\linewidth]{img/2008/fujian_science/q21_fig1.png}
\end{problem}

\begin{solution}
设椭圆半焦距为 \(c\). 因为焦点为 \(F(1,0)\)，所以 \(c=1\)，从而
\[
a^2-b^2=1.
\]
\begin{enumerate}
\item 设短轴的两个三等分点为 \(M\left(0,\frac{b}{3}\right)\) 和 \(N\left(0,-\frac{b}{3}\right)\)，取焦点 \(F(1,0)\). 则 \(MN=\frac{2b}{3}\)，\(FM=FN=\sqrt{1+\frac{b^2}{9}}\).
三角形 \(FMN\) 为正三角形，所以
\[
\frac{2b}{3}=\sqrt{1+\frac{b^2}{9}}.
\]
两边平方得 \(\frac{4b^2}{9}=1+\frac{b^2}{9}\)，即 \(b^2=3\). 由 \(a^2-b^2=1\)，得 \(a^2=4\). 因为 \(a>0\)，所以椭圆方程为
\[
\frac{x^2}{4}+\frac{y^2}{3}=1.
\]

\item 设过点 \(F\) 的直线 \(l\) 的单位方向向量为 \((u,v)\)，其中 \(u^2+v^2=1\)，则直线上的点可表示为 \((1+ut,vt)\). 设它与椭圆交于对应参数 \(t_1,t_2\) 的点 \(A,B\)，并令
\[
K=b^2u^2+a^2v^2=a^2-u^2.
\]
将直线参数式代入椭圆方程，并利用 \(a^2-b^2=1\)，得
\[
Kt^2+2b^2ut-b^4=0.
\]
所以 \(t_1+t_2=-\frac{2b^2u}{K}\)，\(t_1t_2=-\frac{b^4}{K}\).
又因为 \(|OA|^2=1+2ut_1+t_1^2\)，\(|OB|^2=1+2ut_2+t_2^2\)，\(|AB|^2=(t_1-t_2)^2\)，
题设不等式等价于
\[
1+u(t_1+t_2)+t_1t_2<0.
\]
代入根与系数的关系，得到
\[
a^2-b^4-(1+2b^2)u^2<0.
\]
左端关于 \(u^2\) 单调递减，而直线任意转动时 \(u^2\in[0,1]\)，故对所有直线恒成立的充要条件是 \(u=0\) 时也成立，即
\[
a^2-b^4<0.
\]
代入 \(b^2=a^2-1\)，得
\[
(a^2-1)^2>a^2.
\]
令 \(z=a^2\)，且 \(z>1\)，则 \(z^2-3z+1>0\). 该不等式的根为 \(\frac{3-\sqrt{5}}{2}\) 和 \(\frac{3+\sqrt{5}}{2}\)，结合 \(z>1\)，可得
\[
a^2>\frac{3+\sqrt{5}}{2}.
\]
因此 \(a\) 的取值范围为 \(\left(\sqrt{\frac{3+\sqrt{5}}{2}},+\infty\right)\).
\end{enumerate}
\end{solution}

\begin{problem}
已知函数 \( f(x) = \ln (1 + x) - x \).
\begin{enumerate}
\item 求 \( f(x) \) 的单调区间；
\item 记 \( f(x) \) 在区间 \([0, \pi]\)（\(n \in \N^*\)）上的最小值为 \(b_n\)，令 \(a_n = \ln(1 + n) - b_n\).
\begin{enumerate}
    \item 如果对一切 \(n\)，不等式 \( \sqrt{a_n} < \sqrt{a_{n+2}} - \frac{c}{\sqrt{a_{n+2}}} \) 恒成立，求实数 \(c\) 的取值范围；
    \item 求证： \(\frac{a_1}{a_2} + \frac{a_1a_3}{a_2a_4} + \dots + \frac{a_1a_3\cdots a_{2n-1}}{a_2a_4\cdots a_{2n}} < \sqrt{2a_n + 1} -1\).
\end{enumerate}
\end{enumerate}
\end{problem}

\begin{solution}
\begin{enumerate}
\item 函数定义域为 \((-1,+\infty)\)，且
\[
f'(x)=\frac{1}{1+x}-1=-\frac{x}{1+x}.
\]
当 \(-1<x<0\) 时，\(f'(x)>0\)；当 \(x>0\) 时，\(f'(x)<0\). 因此 \(f(x)\) 在 \((-1,0)\) 上单调递增，在 \((0,+\infty)\) 上单调递减.

\item 因为 \(f(x)\) 在 \([0,+\infty)\) 上单调递减，所以对任意 \(n\in\N^*\)，其在 \([0,\pi]\) 上的最小值为
\[
b_n=f(\pi)=\ln(1+\pi)-\pi.
\]
令
\[
d=\pi-\ln(1+\pi)>0,
\]
则
\[
a_n=\ln(1+n)-b_n=d+\ln(1+n)>0.
\]

\begin{enumerate}
\item 因为 \(a_{n+2}>a_n>0\)，原不等式等价于
\[
c<a_{n+2}-\sqrt{a_na_{n+2}}=:h_n.
\]
又
\[
h_n=(a_{n+2}-a_n)\frac{\sqrt{a_{n+2}}}{\sqrt{a_{n+2}}+\sqrt{a_n}},
\]
所以
\[
0<h_n<a_{n+2}-a_n=\ln\frac{n+3}{n+1}.
\]
由 \(\displaystyle\lim_{n\to+\infty}\ln\frac{n+3}{n+1}=0\)，若 \(c>0\)，则对充分大的 \(n\) 有 \(h_n<c\)，原不等式不能对一切 \(n\) 成立. 若 \(c\leqslant0\)，则 \(c<h_n\) 对一切 \(n\in\N^*\) 成立. 因此
\[
c\in(-\infty,0].
\]

\item 按原题面，第二问中的不等式并不成立. 取 \(n=2\)，由
\[
a_1=d+\ln2\approx2.413659,\quad a_2=d+\ln3\approx2.819125,
\]
\[
a_3=d+\ln4\approx3.106807,\quad a_4=d+\ln5\approx3.329950,
\]
左端为
\[
\frac{a_1}{a_2}+\frac{a_1a_3}{a_2a_4}\approx1.654974,
\]
而右端为
\[
\sqrt{2a_2+1}-1\approx1.576480.
\]
于是左端大于右端，所要求证的不等式在 \(n=2\) 时已经不成立. 因此按原题清晰题面，第二问存在排印或版本错误，不能给出其所要求的不等式证明.
\end{enumerate}
\end{enumerate}
\end{solution}
